\section{Conclusion}
\label{sec:conclusion}

\subsection{Future work}

The slicing algorithms presented in \secref{graph} are neutral with respect to any intensional information
associated with the dependence graph, in the sense that the Galois connections they induce are determined only
by the \emph{extension} -- I/O relation -- of the graph, rather than any internal paths. It follows that
transformations that elide internal nodes but preserve I/O connectivity, such as the classical Y-$\Delta$
transform \cite{truemper89}, will produce an extensionally equivalent (but faster) analysis. Although we leave
applications of this observation to future work, we conjecture that it might be used to selectively trade
intensional information for additional performance improvements.

The fact that our work is neutral with respect to intensional information in the dependence graph suggests
the question of what information the internal nodes represent. The choice of dependency relation 
we represent in the graph is the same as that of \citet{perera22}, where selections are pushed through
evaluation according to whether the selected part of a value was consumed during evaluation, either via
pattern matching or the specific implementation provided for a primitive. Any information that could be
associated with internal nodes or subgraphs will be determined by the choice of dependency relation, and so
a promising future direction is to formally consider the choice of dependency relation and the information
it allows us to compute about internal nodes in the graph. For example,  \citet{cheney11b} considers
the use of dependency analysis to compute \textit{provenance} information about elements of a program, such as
it's derivation. Applying this concept to the dependence graph would provide an interesting perspective on
what the internal nodes mean, and whether one can apply meaning at the fine grained level of graph nodes,
or if provenance information in the graph can only be considered at a coarser level.  